\vspace{-1em}
\section{Conclusion and Future Work}
\label{sec:conclusion}

Current embodied systems face significant challenges in human-robot interaction, including rigid working patterns and the inability to handle interruptions. The VITA-E framework addresses these limitations through its innovative dual-model architecture and special token-based control flow, taking an important step toward achieving concurrent and interruptible human-robot interaction. This is enabled by our proposed ``model-as-controller'' paradigm, where the VLM generates its own control tokens to directly control system behavior, tightly coupling high-level reasoning with system execution. While this dual-model architecture provides robustness and an effective solution to interruption, we acknowledge that it comes at the cost of higher computational resource consumption compared to single-model systems. There are also several limitations in our current implementation: further enhancement of model capabilities or verification of framework universality.

Although more work remains, this work opens several exciting avenues for future research. First, our architecture could be extended to handle long-horizon, multi-stage tasks by exploiting the high-level VLM to direct the low-level execution policies step-by-step. Additionally, the framework's inherent support for interruption is well-suited for incorporating real-time human feedback to correct erroneous actions or guide novel behaviors, thereby improving task success rates. Finally, while our current task-switching relies on a safe retraction to a neutral state, we plan to explore methods for achieving smoother and more efficient transitions between tasks.

We envision that this work will inspire future research toward more natural and intelligent human-robot collaboration, ultimately realizing the goal of seamless embodied assistants that can adapt and respond to human needs in dynamic, real-time environments.
